\documentclass[11pt]{article}

\usepackage[a4paper,margin=25mm]{geometry}
\usepackage[T1]{fontenc}
\usepackage[utf8]{inputenc}
\usepackage{lmodern}
\usepackage{microtype}
\usepackage{amsmath,amssymb,mathtools}
\usepackage{amsthm}
\usepackage{array}
\usepackage{booktabs}
\usepackage{enumitem}
\usepackage{tabularx}
\usepackage{xcolor}
\usepackage{hyperref}

\definecolor{artifactblue}{HTML}{1F4E79}
\definecolor{artifactgreen}{HTML}{1E6B45}
\definecolor{artifactred}{HTML}{8B2F2F}

\hypersetup{
  colorlinks=true,
  linkcolor=artifactblue,
  urlcolor=artifactblue,
  pdftitle={A Mathematical Reading of TargetKeygenSamplerCallers.ec},
  pdfsubject={HAETAE sampler-caller schedule refinement}
}

\setlist[itemize]{leftmargin=1.5em,itemsep=0.25em,topsep=0.35em}
\setlist[enumerate]{leftmargin=1.7em,itemsep=0.25em,topsep=0.35em}
\renewcommand{\arraystretch}{1.18}
\setlength{\emergencystretch}{3em}

\theoremstyle{definition}
\newtheorem{definition}{Definition}
\theoremstyle{plain}
\newtheorem{theorem}{Theorem}
\theoremstyle{remark}
\newtheorem*{remark}{Remark}

\newcommand{\code}[1]{\nolinkurl{#1}}
\newcommand{\Z}{\mathbb Z}
\newcommand{\W}{\mathbb W_{64}}
\newcommand{\enc}[1]{\left[#1\right]_{64}}
\newcommand{\uint}[1]{\langle #1\rangle}
\newcommand{\AL}{\mathcal A_{\mathrm L}}
\newcommand{\AS}{\mathcal A_{\mathrm S}}
\newcommand{\Seed}{\mathcal S}
\newcommand{\UL}{\mathcal U_{\mathrm L}}
\newcommand{\US}{\mathcal U_{\mathrm S}}
\newcommand{\ES}{\mathcal E_{\mathrm S}}
\newcommand{\SchedMat}{\mathsf M}
\newcommand{\SchedVec}{\mathsf V}
\newcommand{\SchedEta}{\mathsf E}
\newcommand{\GenMat}{\widehat{\mathsf M}}
\newcommand{\GenVec}{\widehat{\mathsf V}}
\newcommand{\GenEta}{\widehat{\mathsf E}}

\newenvironment{resultbox}
  {\begin{center}\begin{minipage}{0.94\linewidth}\color{artifactgreen}\hrule\vspace{0.6em}\color{black}}
  {\vspace{0.6em}\color{artifactgreen}\hrule\end{minipage}\end{center}}

\newenvironment{boundarybox}
  {\begin{center}\begin{minipage}{0.94\linewidth}\color{artifactred}\hrule\vspace{0.6em}\color{black}}
  {\vspace{0.6em}\color{artifactred}\hrule\end{minipage}\end{center}}

\title{Three Array Schedules in HAETAE\\
\large A mathematical reading of \code{TargetKeygenSamplerCallers.ec}}
\author{HAETAE reference EasyCrypt artifact}
\date{Artifact state checked 2026-07-23}

\begin{document}
\maketitle

\begin{abstract}
The verified program contains three loops that repeatedly call polynomial
sampling subroutines.  The relational theorems check their arithmetic and
ordering against transparent integer schedules.  Separate Hoare theorems in
the same refinement carry exact finite SHAKE decoder prefixes, coefficient
ranges, and aggregate frames through the two actual uniform callers and the
actual eta caller.  These are deterministic schedule and partial-correctness
results, not sampler-distribution theorems.
\end{abstract}

\begin{resultbox}
\textbf{The result in one sentence.}
The generated matrix, vector, and eta callers have the same returned-array
behavior as the explicit schedules in Definition~\ref{def:schedules}; the
proof uses no semantic property of a sampler beyond applying the same fixed
procedure to equal arguments.
\end{resultbox}

\section{The Mathematical Question}

Each caller starts with an array and a seed buffer.  It repeatedly invokes a
lower-level subroutine with five arguments:
\[
  (\text{current array},\ \text{output base},\ \text{seed buffer},\
    \text{seed offset},\ \text{nonce}).
\]
The lower-level subroutine returns a new array, which becomes the input to the
next call.

Jasmin is the language of the implementation.  Here ``generated'' or
``extracted'' means mechanically translated from Jasmin into EasyCrypt, while
the ``authored schedule'' is a short, hand-written EasyCrypt loop used for
comparison.

The verification question is:

\begin{quote}
Can every iteration of the generated word-level loop be paired with the
corresponding integer-indexed iteration, so that both sides pass the same five
arguments to the same subroutine and ultimately return equal arrays?
\end{quote}

There are three callers:
\begin{itemize}
\item a nested loop that constructs a matrix of polynomials;
\item a single loop that constructs a vector using the uniform sampler; and
\item a single loop that constructs a vector using the eta sampler.
\end{itemize}

The proof answers this question affirmatively at the level of the returned
array.  It does not export a separate theorem about call traces.  The authored
schedule is a code-derived normal form of these loops, not an independent
specification of HAETAE from the paper.

\section{A Model with Opaque Sampling Subroutines}

Let
\[
  \W = \Z/2^{64}\Z
\]
be the set of 64-bit words.  For an integer $x$, write $\enc{x}$ for its
residue in $\W$.  For $w\in\W$, write $\uint{w}$ for its canonical
representative in $\{0,\ldots,2^{64}-1\}$.

Let $\AL$ be the set of large arrays, $\AS$ the set of small arrays, and
$\Seed$ the set of seed buffers.  In the extracted model, a large array has
8192 32-bit slots, a small array has 2048 32-bit slots, and a seed buffer has
128 bytes.  A base is an index measured in 32-bit slots, whereas a seed offset
is measured in bytes.

For sets $X$ and $Y$, the notation $X\rightharpoonup Y$ means a function that
may be undefined on some inputs.  We use this notation because a sampling
procedure may fail to return.  We name the three lower-level procedures as
partial maps:
\begin{align*}
  \UL &: \AL\times\W\times\Seed\times\W\times\W
          \rightharpoonup \AL,\\
  \US &: \AS\times\W\times\Seed\times\W\times\W
          \rightharpoonup \AS,\\
  \ES &: \AS\times\W\times\Seed\times\W\times\W
          \rightharpoonup \AS.
\end{align*}
Here $\UL$ is the uniform-sampling procedure for a large output array, $\US$
is its small-array counterpart, and $\ES$ is the eta-sampling procedure.  The
symbol $\rightharpoonup$ is deliberate: the present artifact does not prove
that the rejection loops inside these procedures terminate unconditionally.
Separate leaf pHoare theorems prove probability-one termination under explicit
deterministic progress certificates, but those certificate premises are not
part of the caller-schedule equivalences.  A later parent refinement lifts
bundled certificates through the fixed \(2\times3\) matrix, two-slot vector,
and first-attempt eta callers.

These symbols denote the three fixed procedures extracted from the Jasmin
source; the theorem is not universally quantified over substitute functions.
Nothing in the main refinement theorem describes the fixed maps further, so
they may be viewed as black boxes for purposes of this proof.  In particular,
the proof does not assume that they are uniform samplers, or even that they
modify only the 256 positions starting at the supplied base.

\section{The Three Transparent Schedules}
\label{sec:schedules}

The following definitions express exactly what the authored EasyCrypt schedule
programs do.  All bases and nonces are words in $\W$, so the brackets
$\enc{\cdot}$ include reduction modulo $2^{64}$.

\begin{definition}[Matrix, vector, and eta schedules]
\label{def:schedules}
Let $S$ be a seed buffer.

\paragraph{Matrix schedule.}
Let $A\in\AL$ and $R,C\in\W$, and put $r=\uint{R}$ and $c=\uint{C}$.
Define $A^{(0)}=A$.  For every row $0\le i<r$, set
\[
  A^{(i,0)}=A^{(i)},
\]
and, for $0\le j<c$, define
\[
  A^{(i,j+1)}=
  \UL\!\left(
    A^{(i,j)},
    \enc{256(ic+j)},
    S,
    \enc{0},
    \enc{256i+j}
  \right).
\]
After the row is complete, set $A^{(i+1)}=A^{(i,c)}$.  The matrix schedule is
\[
  \SchedMat(A,S,R,C)=A^{(r)}.
\]

\paragraph{Uniform-vector schedule.}
Let $A\in\AS$ and $K,M\in\W$, and put $k=\uint{K}$ and $m=\uint{M}$.
Starting from $B^{(0)}=A$, define for $0\le i<k$
\[
  B^{(i+1)}=
  \US\!\left(
    B^{(i)},
    \enc{256i},
    S,
    \enc{0},
    \enc{256k+m+i}
  \right).
\]
The vector schedule is
\[
  \SchedVec(A,S,K,M)=B^{(k)}.
\]

\paragraph{Eta-vector schedule.}
Let $A\in\AS$ and $N,T\in\W$, where $N$ is the initial nonce, and put
$t=\uint{T}$.  Starting from $D^{(0)}=A$, define for $0\le i<t$
\[
  D^{(i+1)}=
  \ES\!\left(
    D^{(i)},
    \enc{256i},
    S,
    \enc{32},
    N+\enc{i}
  \right).
\]
The eta schedule is
\[
  \SchedEta(A,S,N,T)=D^{(t)}.
\]
\end{definition}

Each recurrence is itself partial: if one displayed sampler call does not
return, that recurrence and its final schedule value are undefined.

These definitions make the essential arithmetic visible:
\begin{center}
\begin{tabular}{llll}
\toprule
Schedule & Indices & Output base & Nonce \\
\midrule
Matrix & $(i,j)$ & $256(ic+j)$ & $256i+j$ \\
Uniform vector & $i$ & $256i$ & $256k+m+i$ \\
Eta vector & $i$ & $256i$ & $N+i$ in $\W$ \\
\bottomrule
\end{tabular}
\end{center}
The matrix and uniform-vector schedules use seed offset $0$; the eta schedule
uses seed offset $32$.

\section{The Theorem Proved by the Refinement File}

Let $\GenMat$, $\GenVec$, and $\GenEta$ denote the three procedures generated
from the Jasmin source.  They use word-valued counters and update their bases
and nonces incrementally.

\begin{theorem}[Caller-schedule refinement]
\label{thm:main}
As partial transformations of their input arrays,
\begin{align*}
  \GenMat(A,S,R,C) &= \SchedMat(A,S,R,C),\\
  \GenVec(A,S,K,M) &= \SchedVec(A,S,K,M),\\
  \GenEta(A,S,N,T) &= \SchedEta(A,S,N,T).
\end{align*}
The matrix equality holds for $A\in\AL$, $S\in\Seed$, and $R,C\in\W$; the
vector equality for $A\in\AS$, $S\in\Seed$, and $K,M\in\W$; and the eta
equality for $A\in\AS$, $S\in\Seed$, and $N,T\in\W$.  On each such input, the
two sides have the same definedness and, when they return, the same array.  The
theorem does not assert a separate equality of call traces.
\end{theorem}

Thus equality as partial transformations includes whether the two sides
return, without claiming that the samplers terminate on every input (called
``losslessness'' in EasyCrypt).  Certificate-conditioned leaf totality does
not change this unconditional, parameterized caller-equivalence reading.  A
separate parent-prefix refinement proves fixed mode-2 caller totality under
bundled certificates.  If a common black-box call fails to return outside
those premises, both sides have reached that same call from equal arguments.
The EasyCrypt judgments are the formal source of
Theorem~\ref{thm:main}; the notation above is its deterministic mathematical
reading.

\begin{remark}[No hidden size assumptions]
The three equivalence statements do not assume $rc\le32$, $k\le8$, or
$t\le8$.  Their formulas are word equalities and are therefore interpreted
modulo $2^{64}$.  Separate arithmetic lemmas show that the intended HAETAE
parameters fit the array capacities, but those lemmas are not hypotheses of
Theorem~\ref{thm:main}.
\end{remark}

\section{Why the Proof Is an Induction}

The checked relational loop invariants have the mathematical content of three
elementary induction arguments.  Their only fact about a black box is this:
identical five-argument tuples give calls with identical definedness and, when
the calls return, identical outputs.  In EasyCrypt this appears as three
self-equivalence lemmas, one for each of $\UL$, $\US$, and $\ES$.

\subsection{Uniform-vector induction}

At the start of iteration $i$, the induction hypothesis identifies the
generated state with the transparent schedule as follows:
\[
\begin{array}{c|c}
\text{generated word variable} & \text{mathematical value} \\
\hline
\text{counter} & \enc{i} \\
\text{base} & \enc{256i} \\
\text{nonce} & \enc{256k+m+i} \\
\text{seed offset} & \enc{0}.
\end{array}
\]
The two current arrays and the seed buffers are also equal.  Consequently both
sides call $\US$ with the same five arguments and obtain the same next array.
The generated program then adds $1$ to its counter and nonce and adds $256$ to
its base.  These updates turn the displayed formulas for $i$ into the formulas
for $i+1$.  The base case $i=0$ is exactly the generated initialization
\[
  \text{base}=\enc{0},\qquad \text{nonce}=\enc{256k+m}.
\]
The implementation realizes multiplication by $256$ as an eight-bit left
shift, so its source expression is $(K\ll8)+M$.

\subsection{Eta-vector induction}

The eta caller has the same structure, but its initial nonce is an input $N$.
The transparent schedule keeps $N$ fixed, while the generated program mutates
its local nonce.  The induction hypothesis is
\[
  \text{counter}=\enc{i},\qquad
  \text{base}=\enc{256i},\qquad
  \text{current nonce}=N+\enc{i},\qquad
  \text{seed offset}=\enc{32}.
\]
Both sides therefore make the same call to $\ES$.  Incrementing the generated
counter and nonce and advancing the base by $256$ proves the induction step.

\subsection{Nested matrix induction}

The outer induction hypothesis, at the start of row $i$, is
\[
  \text{row counter}=\enc{i},\qquad
  \text{row base}=\enc{ic},\qquad
  \text{seed offset}=\enc{0}.
\]
Inside that row, the inner induction adds
\[
  \text{column counter}=\enc{j}.
\]
The generated program computes
\[
  \text{base}=\enc{256(ic+j)},\qquad
  \text{nonce}=\enc{256i+j},
\]
so both sides call $\UL$ with the same arguments.  Increasing $j$ proves the
inner induction step.  At the end of the row, the generated update
\[
  \text{row base}\leftarrow \text{row base}+C
\]
changes $\enc{ic}$ into $\enc{(i+1)c}$, which proves the outer step.

\subsection{Why the loop guards agree}

The generated loops compare 64-bit words, whereas Definitions~\ref{def:schedules}
use integers.  During a loop, the counter relation is $\text{counter}=\enc{i}$
with $0\le i\le\uint{K}$ (or the corresponding row, column, or count bound).
On this range,
\[
  \enc{i}<_{\mathrm u}K
  \quad\Longleftrightarrow\quad
  i<\uint{K},
\]
where $<_{\mathrm u}$ is unsigned word comparison.  Hence the generated loop
and the transparent loop stop at the same index.

\section{A Concrete Mode-2 Example}

The companion specification contains separate arithmetic lemmas for the
intended mode-2 values $k=2$ and $m=3$.  These lemmas are not used by
Theorem~\ref{thm:main}, and no current EasyCrypt theorem proves that the full
mode-2 key generator supplies these parameters.  Substituting
$R=\enc{2},C=\enc{3}$ in the matrix schedule gives two rows and three columns:

\begin{center}
\begin{tabular}{cccc}
\toprule
$(i,j)$ & Linear position $ic+j$ & Base & Nonce \\
\midrule
$(0,0)$ & $0$ & $0$ & $0$ \\
$(0,1)$ & $1$ & $256$ & $1$ \\
$(0,2)$ & $2$ & $512$ & $2$ \\
$(1,0)$ & $3$ & $768$ & $256$ \\
$(1,1)$ & $4$ & $1024$ & $257$ \\
$(1,2)$ & $5$ & $1280$ & $258$ \\
\bottomrule
\end{tabular}
\end{center}

Substituting $K=\enc{2},M=\enc{3}$ in the uniform-vector schedule gives two
calls:
\begin{center}
\begin{tabular}{ccc}
\toprule
$i$ & Base $256i$ & Nonce $256k+m+i$ \\
\midrule
$0$ & $0$ & $515$ \\
$1$ & $256$ & $516$ \\
\bottomrule
\end{tabular}
\end{center}

For a retry index $q$, a separate arithmetic lemma proves the five eta nonce
values
\[
  5q,\quad 5q+1,\quad 5q+2,\quad 5q+3,\quad 5q+4.
\]
It does \emph{not} prove that the complete key-generation program groups the
first three values into one eta call and the last two into another.  Nor does
Theorem~\ref{thm:main} prove that the complete program supplies $k=2$, $m=3$,
or this retry index.  Those are later composition obligations.

The companion theory also proves the following independent arithmetic facts:
\begin{itemize}
\item if $0\le r,c$, $rc\le32$, $0\le i<r$, and $0\le j<c$, the matrix base
for $(i,j)$ through base$+255$ lies below $8192$;
\item if $0\le q\le8$ and $0\le i<q$, the linear base for $i$ through
base$+255$ lies below $2048$; and
\item the 256-slot regions for distinct nonnegative linear indices are
disjoint.
\end{itemize}
These facts are machine checked, but the three refinement proofs do not invoke
them.

\section{What the Theorem Does Not Prove}

\begin{boundarybox}
Theorem~\ref{thm:main} proves the organization of the caller loops relative to
the same opaque subroutines.  It does not prove that those subroutines have the
cryptographic or memory properties suggested by their names.
\end{boundarybox}

\noindent\begin{tabularx}{\textwidth}{>{\raggedright\arraybackslash}p{0.43\textwidth}X}
\toprule
Question & Status \\
\midrule
Do $\UL$ and $\US$ implement the exact finite SHAKE128 decoder path? &
Proved on return by the separate stream Hoare theorems. \\
Does $\ES$ implement SHAKE256 and the intended eta distribution? &
The exact finite SHAKE256 decoder path is proved on return; its probability distribution is not. \\
Do the sampling loops terminate on every input (are they ``lossless'' in
EasyCrypt)? &
Not unconditionally.  Separate leaf theorems are probability one only under
explicit deterministic progress certificates.  A separate parent refinement
proves totality of the fixed mode-2 callers and proof-only prefix under a
bundle of thirteen such certificates, not for every input. \\
Does a call modify only the intended 256-slot output region? &
Yes in the separate leaf and aggregate caller frame theorems. \\
Are physical pointers separated and memory accesses safe? &
The extraction uses value arrays, so source-level pointer properties are not
represented. \\
Do seed offsets $0$ and $32$ select the paper's intended named seeds? &
The raw offsets and seed-expansion slices are checked, and a separate
proof-only actual-parent prefix composes that expanded buffer into the
sampler calls.  Paper-name identification and full-parent dataflow remain
open. \\
Does the full mode-2 key generator call these procedures with the intended
parameters and retry schedule? &
Its actual 32-file/56-procedure closure is extracted.  A separate proof-only
prefix checks parameters \(k=2,m=3\) and the first five eta nonces, but it is
not \code{_keypair_full_m23} and does not cover rejection retries or later
calls. \\
Does this establish paper-level \code{expandA}/\code{expandS} correctness or a
security-game distribution? &
No. \\
\bottomrule
\end{tabularx}

The extracted model also treats speculative-load-hardening protection
operations as identities and erases source annotations such as
\code{spill}, \code{unspill}, and \code{declassify}.  Therefore this result is
not a constant-time or speculative-execution security proof.

\section{Translation Back to the EasyCrypt Files}

A reader interested only in the mathematics can stop after the previous
section.  The following table locates each mathematical object in the artifact.

\noindent\begin{tabularx}{\textwidth}{>{\raggedright\arraybackslash}p{0.23\textwidth}X}
\toprule
Artifact & Mathematical role \\
\midrule
Generated target & Generated definitions of $\UL$, $\US$, $\ES$, and the three word-level
callers $\GenMat$, $\GenVec$, $\GenEta$. \\
Authored schedule & Definitions of the transparent schedules, word/integer identities, capacity
facts, and concrete mode-2 arithmetic. \\
Authored refinement & The three black-box self-equivalences and the three instances of
Theorem~\ref{thm:main}. \\
\bottomrule
\end{tabularx}

The corresponding paths are:
\begin{center}
\footnotesize
\code{easycrypt/extract/keygen-sampler-callers/KeygenSamplerCallersTarget.ec}\\
\code{easycrypt/spec/KeygenSamplerCallersSpec.ec}\\
\code{easycrypt/refinement/TargetKeygenSamplerCallers.ec}
\end{center}

The theorem-name correspondence is:
\begin{center}
\small
\begin{tabularx}{0.96\textwidth}{>{\raggedright\arraybackslash}p{0.24\textwidth}X}
\toprule
Mathematical statement & EasyCrypt lemma \\
\midrule
$\UL$ agrees with itself on equal inputs &
\code{uniform8192_leaf_self_equiv} \\
$\US$ agrees with itself on equal inputs &
\code{uniform2048_leaf_self_equiv} \\
$\ES$ agrees with itself on equal inputs &
\code{eta2048_leaf_self_equiv} \\
$\GenVec=\SchedVec$ & \code{expand_vecA_relative_orchestration_equiv} \\
$\GenEta=\SchedEta$ & \code{expand_eta_relative_orchestration_equiv} \\
$\GenMat=\SchedMat$ & \code{expand_matA_relative_orchestration_equiv} \\
\bottomrule
\end{tabularx}
\end{center}

An EasyCrypt judgment of the form
\begin{center}
\code{equiv [generated ~ schedule : equal inputs ==> equal result]}
\end{center}
should be read here as equality of the two returned-array behaviors.  The
formal postcondition mentions only the result.  The alignment of calls is the
inductive proof method, not a separately exposed trace statement.

\section{Replaying the Machine-Checked Evidence}

From the \code{haetae-ref-easycrypt/} directory, run
\begin{center}
\code{./scripts/verify-keygen-sampler-callers-proof.sh}
\end{center}
The script starts a Why3 server, checks the pinned sources, rejects extraction
drift, compiles the generated target and the two authored theories with
\code{-no-eco}, and scans the authored files for proof holes and axiom
declarations.

The current unified extraction selects the three caller roots and
\code{_kp_expand_seedbuf}; its exact closure contains 25 generated files and
31 procedures.  The proof manifest contains the generated target plus six
specifications and six refinements.  The separate actual mode-2 parent gate
checks a 32-file, 56-procedure closure.  Its dedicated 17-entry proof gate
also transfers selected contracts to the parent target and checks the
proof-only first-attempt sampler prefix.  Proof-hole and axiom scans do not
establish that every imported dependency is axiom free, and extraction
agreement does not prove the extractor correct.

Build this guide with
\begin{quote}
\begin{verbatim}
cd haetae-ref-easycrypt/latex
make
\end{verbatim}
\end{quote}
The PDF is written to
\code{build/TargetKeygenSamplerCallers.pdf}.

\section*{Conclusion}

The relational core proves that word-level counters, bases, and nonces describe
the same three recurrences as transparent integer counters.  Separate Hoare
proofs now supply exact finite decoder, range, and frame contracts through the
actual callers, and the standalone seed-XOF theorem identifies the 128-byte
expansion and slices on return.  The actual leaves also have
certificate-conditioned probability-one termination theorems.  Universal
certificates, full-parent and retry semantics, source-pointer properties,
sampler distributions, and full key-generation composition remain the next
boundary; a separate proof-only actual-parent prefix now crosses the local
dataflow cut through the first eta attempt.

\end{document}
